\setchapterpreamble[u]{\margintoc}
\chapter{Solving \texorpdfstring{\(A\vect{x}=\vect{b}\)}{Ax=b}: Elimination and
\texorpdfstring{\(A=LU\)}{A=LU}}
\labch{elimination}

\sources{Strang, \emph{Introduction to Linear Algebra}, 6th ed.\ \cite{strang2023ila},
§1.4 and §2.2--2.7; MIT 18.06 \cite{ocw1806}, Lectures~2--5.}

\refch{vectors-and-column-space} explained what a solution \emph{is}. This
chapter is about how to find one without guessing, and about what the procedure
leaves behind: a factorisation of \(A\) that is worth more than the solution it
was invented to produce.

\section{Elimination on an example}
\labsec{elimination-example}

Take
\[
	A=\begin{bmatrix}1&2&1\\3&8&1\\0&4&1\end{bmatrix},
	\qquad
	\vect{b}=\begin{bmatrix}2\\12\\2\end{bmatrix},
\]
and work with the augmented array \(\begin{bmatrix}A&\vect{b}\end{bmatrix}\) so
that the right-hand side receives every operation the left-hand side does.

The first \emph{pivot} is the entry \(a_{11}=1\). Subtracting \(3\times\)row~1
from row~2 clears the entry below it, and the \((3,1)\) entry is already zero, so
nothing is needed there:
\[
	\left[\begin{array}{ccc|c}
		1&2&1&2\\ 3&8&1&12\\ 0&4&1&2
	\end{array}\right]
	\longrightarrow
	\left[\begin{array}{ccc|c}
		1&2&1&2\\ 0&2&-2&6\\ 0&4&1&2
	\end{array}\right].
\]

\marginnote{The multiplier \(3\) is the number of row~1 subtracted from row~2. It
is worth recording rather than discarding: \refsec{lu} assembles all the
multipliers into a matrix.}

The second pivot is the \(2\) now sitting in position \((2,2)\). Subtracting
twice row~2 from row~3 clears below it:
\[
	\longrightarrow
	\left[\begin{array}{ccc|c}
		1&2&1&2\\ 0&2&-2&6\\ 0&0&5&-10
	\end{array}\right]
	=\left[\begin{array}{c|c}U&\vect{c}\end{array}\right].
\]

The left block \(U\) is \emph{upper triangular}, the pivots are \(1,2,5\), and
the transformed right-hand side is \(\vect{c}=[2,6,-10]\tp\).

\begin{definition}
A \emph{pivot} is the first non-zero entry in a row at the moment that row is
used to clear the entries beneath it. The positions of the pivots, not their
values, determine the structure of the solution set.
\end{definition}

\section{Back-substitution}
\labsec{back-substitution}

Triangular systems are solved from the bottom up. Reading \(U\vect{x}=\vect{c}\)
as equations,
\begin{align*}
	x+2y+z &= 2,\\
	2y-2z  &= 6,\\
	5z     &= -10,
\end{align*}
the last equation gives \(z=-2\). Substituting into the second,
\(2y-2(-2)=6\), so \(2y=2\) and \(y=1\). Substituting both into the first,
\(x+2(1)+(-2)=2\), so \(x=2\). Hence
\[
	\vect{x}=\begin{bmatrix}2\\1\\-2\end{bmatrix},
\]
and substituting into the original \(A\vect{x}\) returns
\([2,\,6+8-2,\,4-2]\tp=[2,12,2]\tp=\vect{b}\), as required.

\marginnote{Always substitute back into the \emph{original} matrix, not into
\(U\). An arithmetic slip made during elimination is invisible to a check
performed on the eliminated system.}

\section{Elimination matrices}
\labsec{elimination-matrices}

Each elimination step is itself a matrix multiplication. Subtracting three
copies of row~1 from row~2 is achieved by left multiplication with
\[
	E_{21}=\begin{bmatrix}1&0&0\\-3&1&0\\0&0&1\end{bmatrix},
\]
which differs from the identity in exactly one off-diagonal entry. Indeed
\(E_{21}A\) reproduces the first array computed above. The second step is
\[
	E_{32}=\begin{bmatrix}1&0&0\\0&1&0\\0&-2&1\end{bmatrix},
	\qquad
	E_{32}\bigl(E_{21}A\bigr)=U .
\]

Because matrix multiplication is associative, the two steps can be combined
before touching \(A\):
\[
	\bigl(E_{32}E_{21}\bigr)A=U .
\]

\begin{remark}
Associativity is doing real work here. It is what allows a sequence of row
operations to be replaced by a single matrix, and it is why the whole
elimination process can be recorded as a factorisation instead of as a script of
instructions.
\end{remark}

Undoing a step is as easy as performing it: if \(E_{21}\) subtracts three copies
of row~1 from row~2, then adding them back inverts it,
\[
	E_{21}^{-1}=\begin{bmatrix}1&0&0\\3&1&0\\0&0&1\end{bmatrix},
	\qquad
	E_{21}^{-1}E_{21}=I .
\]

\section{Matrix multiplication, four ways}
\labsec{four-ways}

Since elimination is multiplication, it is worth collecting the ways a product
\(C=AB\) with \(A\in\R^{m\times n}\) and \(B\in\R^{n\times p}\) can be read. All
four produce the same \(C\); each makes a different fact obvious.

\begin{description}
	\item[Entry by entry] \(c_{ij}=\sum_{k=1}^{n}a_{ik}b_{kj}\), the dot product
	of row \(i\) of \(A\) with column \(j\) of \(B\).
	\item[Column by column] column \(j\) of \(C\) is \(A\vect{b}_j\). Hence
	\emph{every column of \(C\) is a combination of the columns of \(A\)}, and
	\(\Col(AB)\subseteq\Col(A)\).
	\item[Row by row] row \(i\) of \(C\) is \(\vect{a}_i\tp B\), a combination of
	the rows of \(B\). Hence the row space of \(AB\) sits inside the row space of
	\(B\).
	\item[Column times row] \(C=\sum_{k=1}^{n}\vect{a}_k\vect{b}_k\tp\), a sum of
	\(n\) rank-one matrices, where \(\vect{a}_k\) is column \(k\) of \(A\) and
	\(\vect{b}_k\tp\) is row \(k\) of \(B\).
\end{description}

\marginnote{The fourth reading is the one that generalises furthest. The
singular value decomposition in \refch{svd} is precisely a column-times-row
expansion in which the terms are ordered by importance.}

\begin{example}
With
\[
	A=\begin{bmatrix}2&7\\3&8\\4&9\end{bmatrix},
	\qquad
	B=\begin{bmatrix}1&6\\0&0\end{bmatrix},
\]
the column-times-row reading gives
\[
	AB=\begin{bmatrix}2\\3\\4\end{bmatrix}\begin{bmatrix}1&6\end{bmatrix}
	+\begin{bmatrix}7\\8\\9\end{bmatrix}\begin{bmatrix}0&0\end{bmatrix}
	=\begin{bmatrix}2&12\\3&18\\4&24\end{bmatrix}.
\]
The second term vanishes because the second row of \(B\) is zero, which the
entry-by-entry rule would have obscured behind six separate multiplications.
\end{example}

A fifth reading, block multiplication, follows from the same algebra: if \(A\)
and \(B\) are cut into conforming blocks then
\[
	\begin{bmatrix}A_1&A_2\\A_3&A_4\end{bmatrix}
	\begin{bmatrix}B_1&B_2\\B_3&B_4\end{bmatrix}
	=\begin{bmatrix}A_1B_1+A_2B_3 & A_1B_2+A_2B_4\\
	                A_3B_1+A_4B_3 & A_3B_2+A_4B_4\end{bmatrix},
\]
provided the block shapes match. Every practical implementation of matrix
multiplication --- and every batched computation in a neural network --- relies
on this.

\section{Row exchanges and permutation matrices}
\labsec{permutations}

A zero in the pivot position is not fatal on its own. If a row further down has
a non-zero entry in that column, exchanging the two rows restores a usable
pivot. The exchange is again a matrix: for two rows,
\[
	P=\begin{bmatrix}0&1\\1&0\end{bmatrix},
	\qquad
	P\begin{bmatrix}a&b\\c&d\end{bmatrix}=\begin{bmatrix}c&d\\a&b\end{bmatrix}.
\]

\begin{remark}
Multiplying on the \emph{left} rearranges rows; multiplying on the \emph{right}
rearranges columns. The same \(P\) gives
\[
	\begin{bmatrix}a&b\\c&d\end{bmatrix}P=\begin{bmatrix}b&a\\d&c\end{bmatrix}.
\]
Side matters, and the habit of checking which side an operation acts on prevents
a large fraction of avoidable errors.
\end{remark}

\begin{definition}
A \emph{permutation matrix} is obtained from \(I\) by reordering its rows. Every
permutation matrix satisfies \(P^{-1}=P\tp\), because the reordering is undone
by the reverse reordering.
\end{definition}

With exchanges allowed, elimination on an invertible matrix always succeeds, and
the factorisation of \refsec{lu} becomes \(PA=LU\).

\section{When elimination fails}
\labsec{elimination-failure}

Failure occurs only when a pivot position holds zero \emph{and} every entry
below it in that column is zero as well: there is nothing left to exchange with.
A schematic third stage of the form
\[
	\begin{bmatrix}\ast&\ast&\ast\\ 0&\ast&\ast\\ 0&0&0\end{bmatrix}
\]
has only two pivots for three columns. The matrix is then singular, and by
\refsubsec{reachability} some targets are unreachable while others are reached in
infinitely many ways.

\begin{remark}
It is worth separating two situations that look alike. A zero \emph{encountered}
in a pivot position is a bookkeeping accident, repaired by a row exchange. A
zero that \emph{survives} every available exchange is a statement about the
matrix: its columns are dependent, and no algorithm can repair that.
\end{remark}

\section{Inverses and Gauss--Jordan}
\labsec{inverses}

For a square matrix, \(A^{-1}\) is the matrix with \(A^{-1}A=I=AA^{-1}\). It
exists exactly when \(A\) is non-singular. The quickest certificate of
\emph{non}-existence is a non-zero vector in the nullspace.

\begin{proposition}
If \(A\vect{x}=\vect{0}\) for some \(\vect{x}\neq\vect{0}\), then \(A\) has no
inverse.
\end{proposition}

\begin{proof}
Suppose \(A^{-1}\) existed. Multiplying \(A\vect{x}=\vect{0}\) on the left gives
\(\vect{x}=A^{-1}\vect{0}=\vect{0}\), contradicting \(\vect{x}\neq\vect{0}\).
\end{proof}

\begin{example}
For \(A=\begin{bmatrix}1&3\\2&6\end{bmatrix}\) the vector
\(\vect{x}=\begin{bmatrix}3&-1\end{bmatrix}\tp\) gives
\[
	\begin{bmatrix}1&3\\2&6\end{bmatrix}\begin{bmatrix}3\\-1\end{bmatrix}
	=\begin{bmatrix}0\\0\end{bmatrix},
\]
so this \(A\) is singular. The second column is three times the first, which is
the same fact stated in the column picture.
\end{example}

To \emph{find} an inverse, note that column \(j\) of \(A^{-1}\) is the solution
of \(A\vect{x}=\vect{e}_j\), where \(\vect{e}_j\) is column \(j\) of \(I\). The
\(n\) systems share the same left-hand side, so they can be eliminated together.
Gauss--Jordan places \(A\) and \(I\) side by side and reduces until the left
block is the identity:
\[
	\begin{bmatrix}A&I\end{bmatrix}\longrightarrow\begin{bmatrix}I&A^{-1}\end{bmatrix}.
\]

\begin{example}
Take \(A=\begin{bmatrix}1&3\\2&7\end{bmatrix}\). Subtracting \(2\times\)row~1
from row~2, then \(3\times\)the new row~2 from row~1:
\[
	\left[\begin{array}{cc|cc}1&3&1&0\\2&7&0&1\end{array}\right]
	\rightarrow
	\left[\begin{array}{cc|cc}1&3&1&0\\0&1&-2&1\end{array}\right]
	\rightarrow
	\left[\begin{array}{cc|cc}1&0&7&-3\\0&1&-2&1\end{array}\right].
\]
So \(A^{-1}=\begin{bmatrix}7&-3\\-2&1\end{bmatrix}\), and multiplying out
confirms \(AA^{-1}=I\).
\end{example}

\marginnote{Why the method works: the row operations amount to left
multiplication by some matrix \(E\) with \(EA=I\), and that identity says
\(E=A^{-1}\). The right-hand block simply records \(E\) by tracking what the
same operations do to \(I\).}

Two facts about inverses are used constantly. First, the inverse of a product
reverses the order,
\[
	(AB)^{-1}=B^{-1}A^{-1},
\]
which is checked by multiplying: \(AB\,B^{-1}A^{-1}=AIA^{-1}=I\). Second,
inversion commutes with transposition, \((A\tp)^{-1}=(A^{-1})\tp\).

\begin{remark}
In numerical work one almost never forms \(A^{-1}\) explicitly. Solving
\(A\vect{x}=\vect{b}\) by elimination costs about a third as much as inverting
\(A\) and is better behaved in floating point. The inverse is a tool for
reasoning, not usually for computing.
\end{remark}

\section{The factorisation \texorpdfstring{\(A=LU\)}{A=LU}}
\labsec{lu}

Return to \(\bigl(E_{32}E_{21}\bigr)A=U\) and move the elimination matrices to
the other side:
\[
	A=E_{21}^{-1}E_{32}^{-1}\,U=LU .
\]
The remarkable part is what \(L\) contains. Multiplying the two inverses,
\[
	L=\begin{bmatrix}1&0&0\\3&1&0\\0&0&1\end{bmatrix}
	  \begin{bmatrix}1&0&0\\0&1&0\\0&2&1\end{bmatrix}
	 =\begin{bmatrix}1&0&0\\3&1&0\\0&2&1\end{bmatrix},
\]
whose subdiagonal entries are exactly the multipliers \(3\) and \(2\) used
during elimination --- in their original positions, with no interference between
them. Verifying,
\[
	LU=\begin{bmatrix}1&0&0\\3&1&0\\0&2&1\end{bmatrix}
	   \begin{bmatrix}1&2&1\\0&2&-2\\0&0&5\end{bmatrix}
	 =\begin{bmatrix}1&2&1\\3&8&1\\0&4&1\end{bmatrix}=A .
\]

\begin{theorem}
If elimination on \(A\) needs no row exchanges, then \(A=LU\) where \(U\) is the
upper triangular matrix produced by elimination and \(L\) is lower triangular
with ones on the diagonal and the multiplier \(\ell_{ij}\) in position
\((i,j)\). With row exchanges, the same statement holds for a permuted matrix:
\(PA=LU\).
\end{theorem}

\marginnote{The multipliers land in \(L\) unchanged only because the inverses are
multiplied in the reverse order of the eliminations. Multiplying them in the
forward order produces cross terms and destroys the pattern.}

Once \(A=LU\) is available, every subsequent right-hand side is cheap: solve
\(L\vect{c}=\vect{b}\) downwards, then \(U\vect{x}=\vect{c}\) upwards. Both are
triangular, so the expensive part of the work is done once and reused. A
symmetric variant splits the pivots out as \(A=LDU\) with \(D\) diagonal, and for
symmetric positive definite matrices it specialises to \(A=LL\tp\); positive
definiteness is taken up in \refch{eigenvalues}.

\section{Transposes and symmetry}
\labsec{transposes}

The transpose exchanges rows and columns, \((A\tp)_{ij}=a_{ji}\). Its two rules
are
\[
	(AB)\tp=B\tp A\tp,
	\qquad
	(A\tp)\tp=A .
\]
A matrix is \emph{symmetric} when \(A\tp=A\). Symmetric matrices are not a
curiosity: for any \(A\) whatsoever, both \(A\tp A\) and \(AA\tp\) are
symmetric, since
\[
	(A\tp A)\tp=A\tp(A\tp)\tp=A\tp A .
\]

\marginnote{\(A\tp A\) is the matrix behind least squares in
\refch{orthogonality} and behind the singular value decomposition in
\refch{svd}. Almost every time a rectangular matrix has to be turned into a
square one, this is how it is done.}

\section{The cost of elimination}
\labsec{cost}

Clearing the first column of an \(n\times n\) matrix touches roughly \(n^{2}\)
entries, the second about \((n-1)^{2}\), and so on, giving
\[
	n^{2}+(n-1)^{2}+\cdots+1\approx\tfrac{1}{3}n^{3}
\]
operations for the factorisation, while each back-substitution costs about
\(n^{2}\). The cubic term is why reusing \(A=LU\) across many right-hand
sides matters, and why the size of a linear solve grows uncomfortably fast ---
the practical reason that iterative and gradient-based methods dominate once
\(n\) reaches the scale of a machine learning problem.
